En los últimos años, la adopción de herramientas de inteligencia artificial (IA) generativa en la industria del desarrollo de videojuegos ha experimentado un crecimiento acelerado. Sin embargo, esta transición tecnológica desencadenado en una reacción negativo por parte de la comunidad de usuarios, derivando en casos donde las empresas se han visto forzadas a retirar componentes desarrollados con IA para mitigar las criticas negativas recibidas por estos elementos generados artificialmente. A pesar de la relevancia de este fenómeno, la literatura científica actual presenta un vacío respecto a las causas de este rechazo: no existe un consenso empírico sobre si el rechazo proviene de un sesgo actitudinal previo del jugador ante el contenido con IA, o si obedece a limitaciones técnicas y deficiencias tangibles ne la calidad del código y la experiencia interactiva resultante.

Para abordar esta problemática y profundizar en el origen de dicha percepción, esta investigación propone un diseño experimental comparativo mediante el desarrollo de un videojuego de escala reducida en dos variantes: una versión implementada con código y visuales asistidos por modelos de IA generatuva y otra programada y con visuales mediante la utilización de métodos tradicionales. 

La metodología se realizará en el ámbito académico con una muestra representativa de al menos 60 estudiantes de la Universidad del Valle de Guatemala. Los participantes serán evaluados mediante un experimento A/B: un grupo experimental operará bajo total transparencia (informado previamente sobre el proceso de desarrollo de cada versión) y un grupo de control operará bajo una condición ciega (sin información previa). Tras interactuar con ambos prototipos, los sujetos responderán a un instrumento psicométrico validado con escala Likert, diseñado para medir dimensiones clave de la experiencia de usuario (UX), apartado visual (sprites), detección de anomalías y la satisfacción percibida.

La relevancia de este trabajo radica en su capacidad para ofrecer evidencia empírica cuantitativa a la industria del videojuego y a la comunidad académica. Desde la perspectiva práctica, los resultados permitirán a los desarrolladores determinar si el uso de IA para determinados aspectos compromete la experiencia de usuario o si el rechazo del mercado responde primordialmente a un sego de percepción. Con ello, la investigación aportará criterios fundamentados para decidir si la implementación de algoritmos generativos debe optimizarse mediantes nuevos estándares de calidad, replantearse en su flujo de producción o restringirse en módulos específicos.